% ch00-front.tex — בקרת מסמך, מונחים, תוכן עניינים ורשימות
\chapter*{בקרת מסמך ואישור}
\addcontentsline{toc}{chapter}{בקרת מסמך ואישור}
\markboth{בקרת מסמך ואישור}{}

הדוח המכין מגדיר מראש מה ייבנה, איך ייבדק ומה ייחשב הצלחה. אחרי שהמרצה מאשר אותו, נבדוק את הפרויקט מול מה שכתוב בו, וכל שינוי ייעשה רק בבקשת שינוי (\en{Change Request}) כתובה ומאושרת, לפי סעיף~\ref{sec:change}. כל קובצי המקור, הפרומפטים ודוחות סוכני הבדיקה נמצאים במאגר:
\begin{center}\repolink\end{center}

\begin{table}[!htb]\centering
\caption{היסטוריית גרסאות}\label{tab:versions}
\begin{tabularx}{\linewidth}{@{}l l Y@{}}
\toprule
\colhead{גרסה} & \colhead{תאריך} & \colhead{תיאור השינוי} \\
\midrule
\num{1.0} & \num{10.11.2026} & טיוטה מלאה ראשונה \\
\num{1.1} & \num{10.11.2026} & תקציר בשיטת ״השורה התחתונה קודם״, יתרונות מבדלים, הנחות יסוד, מסמך ממשקים, בסיס אומדן ורשימת פעולות פתוחות \\
\num{2.0} & \num{21.11.2026} & מעבר ל-\en{LuaLaTeX} ולתבנית \code{my-template.cls}, פרויקט מסודר (פרק בכל קובץ), מקורות ב-\en{biber}, עדכון פרק המתחרים, וסבב סוכני בדיקה \\
\num{2.1} & \num{22.11.2026} & אחרי סוכן הציון הסופי: מקורות ראשוניים, סעיף אכיפה לתקן (\ref{sec:enforce}), השוק הישראלי (\ref{sec:israel}), \en{AP Verify} בטבלאות המתחרים, קישור למאגר, תבנית בגרסה \num{1.3} \\
\bottomrule
\end{tabularx}
\end{table}

\begin{table}[!htb]\centering
\caption{אישור הדוח}\label{tab:sign}
\begin{tabularx}{\linewidth}{@{}l Y l l@{}}
\toprule
\colhead{תפקיד} & \colhead{שם} & \colhead{חתימה} & \colhead{תאריך} \\
\midrule
מגיש & גיא יערי & \rule{2.6cm}{0.4pt} & \rule{1.8cm}{0.4pt} \\
מגיש & פז ניסים & \rule{2.6cm}{0.4pt} & \rule{1.8cm}{0.4pt} \\
מרצה (מאשר) & ד״ר יורם סגל & \rule{2.6cm}{0.4pt} & \rule{1.8cm}{0.4pt} \\
\bottomrule
\end{tabularx}
\end{table}

\section*{רשימת קיצורים ומונחים}
\begin{table}[!htb]\centering
\begin{tabularx}{\linewidth}{@{}l Y@{}}
\toprule
\colhead{מונח} & \colhead{פירוש} \\
\midrule
\en{AI} (בינה מלאכותית) & מערכות לומדות; בדוח זה בעיקר מחוללי תמונות (\en{Diffusion}, \en{GAN}) \\
\en{C2PA} & \en{Coalition for Content Provenance and Authenticity}: תקן פתוח לחתימה קריפטוגרפית על מקור התוכן \\
\en{Content Credentials} & ״תעודת מקור״ שמוטמעת בקובץ לפי \en{C2PA}: מי צילם, באיזה מכשיר, אילו עריכות בוצעו \\
\en{CMS} & מערכת ניהול התוכן של אתר החדשות \\
\en{ROC} & עקומה שמציגה את שיעור הגילוי (\en{TPR}) מול שיעור התרעות השווא (\en{FPR}) \\
\en{MFA} & \en{Made For Advertising}: אתרים שנבנו כדי לגרוף הכנסות מפרסום פרוגרמטי \\
\en{TRUTH-IMG} & התקן האתי שמוצע בפרויקט זה לטיפול בתמונות במערכות חדשות \\
טראפיק & נפח הכניסות לאתר או לפוסט; המדד שעליו נשען מודל ההכנסות מפרסום \\
\bottomrule
\end{tabularx}
\end{table}

\clearpage
{\small\setlength{\parskip}{0pt}\setstretch{1.0}\tableofcontents}

\clearpage
{\small\setlength{\parskip}{0pt}\listoffigures\listoftables}
